% Generated draft: requires mathematical and source review.
\documentclass[11pt]{article}
\usepackage[margin=25mm]{geometry}
\usepackage{fontspec}
\setmainfont{Latin Modern Roman}
\usepackage{amsmath,amssymb,mathtools,unicode-math}
\setmathfont{Latin Modern Math}
\usepackage{enumitem,needspace,xurl,hyperref}
\setlength{\parindent}{0pt}
\setlength{\parskip}{5pt}
\setlength{\emergencystretch}{3em}
\newcommand{\N}{\mathbb{N}}
\newcommand{\Z}{\mathbb{Z}}
\newcommand{\Q}{\mathbb{Q}}
\newcommand{\R}{\mathbb{R}}
\newcommand{\C}{\mathbb{C}}
\newcommand{\F}{\mathbb{F}}
\newcommand{\A}{\mathbb{A}}
\newcommand{\PP}{\mathbb{P}}
\newcommand{\E}{\mathbb{E}}
\newcommand{\eps}{\varepsilon}
\newcommand{\dd}{\,\mathrm{d}}
\begin{document}
\section*{UnsolvedMath ID: 3800007}
\textbf{Source problem code:} catalogue entry
\subsection{Definitions and mathematical statement}
A usual pseudocircle is a simple closed Jordan curve in the plane; a family of pseudocircles has at most two intersection points for each pair of distinct curves. For an arrangement $\mathcal A$ of $n$ such curves, let $t(\mathcal A)$ count unordered tangent pairs and let $d(\mathcal A)$ count digon cells. A cell is a connected component of the complement of the curves. A digon is a cell whose boundary consists of two arrangement edges, each an arc of a curve between arrangement vertices.

For any fixed admissibility conventions, define
\[
T_{\mathrm{ps}}(n)=\max_{\mathcal A}t(\mathcal A),\qquad
D_{\mathrm{ps}}(n)=\max_{\mathcal A}d(\mathcal A),
\]
where the maxima range over the permitted general-position arrangements of $n$ pseudocircles. Determine both extremal functions. Define $T_{\mathrm{circ}}(n)$ and $D_{\mathrm{circ}}(n)$ in the same way with every curve required to be a genuine Euclidean circle. Determine whether their worst-case bounds coincide; an exact interpretation asks whether $T_{\mathrm{circ}}(n)=D_{\mathrm{circ}}(n)$, while an asymptotic interpretation asks for equality of their sharp growth bounds.

The source does not give the admissibility conventions needed to fix these functions. In particular, ``general position'' must permit the tangencies being counted; a convention requiring only transverse crossings would exclude them. A usual simple tangency is a single common point at which two curves touch without crossing, but the treatment of pairs with more than one contact is unstated. The source also does not say whether the unbounded face is counted as a digon, or whether exact equality or equality of asymptotic bounds is intended.
\subsection{Short English statement}
Find the largest possible numbers of tangent pairs and two-sided cells among $n$ pseudocircles. For ordinary circles, compare the two sharp worst-case bounds. The allowable tangencies and the digon-counting conventions need clarification.
\subsection{Sources}
\begin{itemize}
\item Jeff Erickson, \emph{Open Problems in Computational Geometry}, item ``Tangent pairs of pseudocircles,'' \url{https://jeffe.cs.illinois.edu/open/}; accessed 2026-07-29, HTML extraction. The source warns that its pages have not been systematically updated since 2001.
\item Local registry \texttt{data/unsolvedmath/problems.json}, permanent UnsolvedMath ID $3800007$, source code catalogue entry, \emph{Tangent pairs of pseudocircles}. The proposer field is ``Jeff Erickson and cited proposers.'' Dataset attribution: Ulam AI and the UnsolvedMath Contributors.
\item The supplied research summary mentions bounds for tangencies and digons but supplies no individual publication or numerical bound.
\end{itemize}
\end{document}
